\section{Limitaciones, conexiones y futuro: cuál es el costo de la dataset attention}

La capacidad expresiva de NPT proviene de hacer que los datapoints sean visibles entre sí, y ese mismo diseño crea también el mayor cuello de botella del sistema. Si el lote contiene $b$ datapoints y la hidden dimension de cada fila, después del flatten, es $h=de$, el ABD attention map tiene tamaño $b\times b$, y el cómputo principal puede escribirse de forma aproximada como
\[
\operatorname{Cost}_{\mathrm{ABD}}=O(b^2h),
\qquad
\operatorname{Memory}_{\mathrm{attn}}=O(b^2).
\]

\begin{itemize}
  \item $b$: número de datapoints que participan efectivamente en una relational lookup;
  \item $h$: hidden width de cada row;
  \item aumentar $b$ amplía la cobertura de context candidato, pero con un quadratic cost;
  \item el mini-batch aleatorio reduce el costo, pero puede omitir neighbors potencialmente clave que quedan fuera del lote.
\end{itemize}

\lecturefigure{slide-28-summary-future-work.jpg}{Síntesis de NPT: experimentos, limitación de quadratic scaling y future work.}{Video oficial de Stanford Online, 01:00:42--01:05:36.}

\begin{knowledgebox}{Lectura de la figura: conclusiones y preguntas abiertas por separado}
A la izquierda están los resultados que esta clase ya mostró: dataset-as-input, datapoint self-attention, tabular/image benchmarks e intervention evidence. A la derecha están las limitaciones aún sin resolver: el scaling y tareas nuevas como multi-task, few-shot, semi-supervised y domain adaptation. El future work no es una lista de capacidades actuales.
\end{knowledgebox}

\subsection{Conexión con GNN}

Si cada datapoint se considera un node, ABD equivale a realizar learned message passing sobre un fully connected graph. La diferencia es que muchas GNN reciben sparse edges definidos de antemano, mientras que NPT aprende las relaciones de forma suave a partir del attention score; sin embargo, los graph models modernos también pueden aprender edges, de modo que la frontera entre ambos no es absoluta.

\teachervoice{En la sesión de Q\&A, el expositor acepta la analogía con el fully connected graph y señala que NPT puede entenderse como discovering relational structure. Neil Band menciona Neural Relational Inference: cuando los edges son desconocidos, tanto attention como message passing aprenden interacciones latentes.}

\begin{warningbox}{Un fully connected graph no implica que las relaciones se hayan descubierto correctamente}
Permitir la comunicación entre cualquier par de puntos solo define un hypothesis space. El modelo todavía puede aprender un shortcut, una spurious similarity o una attention casi uniform; para determinar si las relaciones son útiles se necesitan corruption, intervention, OOD evaluation y sparse retrieval escalable.
\end{warningbox}

\subsection{El valor distinto de small data y large data}

La ventaja de small data es que el conjunto de datos completo cabe en un solo forward, de modo que el modelo obtiene un context casi completo; large data, en cambio, nos obliga a introducir retrieval, sparsity, kernel approximation o learned representatives. Aunque estas opciones sacrifican la attention exacta con conexión total, pueden vincular NPT con el retrieval-augmented learning moderno: los parámetros se encargan de aprender las reglas de recuperación y transformación, y los datapoints externos se encargan de almacenar hechos actualizables.

\teachervoice{Al final, el expositor plantea una pregunta incisiva: ¿cuántos parámetros de los grandes parametric models se dedican en realidad a «almacenar datos»? Si una lookup explícita puede asumir parte de la memoria, los parámetros quizá podrían concentrarse más en aprender cómo recuperar, combinar y razonar. Esta pregunta cobró después todavía más importancia en los retrieval-augmented systems.}

\begin{knowledgebox}{Hoja de ruta de scaling}
\begin{tabularx}{\textwidth}{L{0.24\textwidth}>{\raggedright\arraybackslash}X >{\raggedright\arraybackslash}X}
\toprule
Línea & Beneficio & Riesgo nuevo \\
\midrule
Random mini-batch & Implementación sencilla y entrenamiento estable & Un datapoint clave puede no estar en el lote; el context es aleatorio. \\
Sparse / local attention & Reduce $b^2$ a una cantidad pequeña de edges & Requiere definir o aprender un candidate graph. \\
ANN / learned retrieval & Primero recupera las top-$k$ rows y luego aplica NPT reasoning & El recall del retriever se vuelve el cuello de botella previo. \\
Representative points & Comprime el conjunto de datos con prototypes / inducing points & La compresión puede perder muestras raras y detalles locales. \\
Kernel / low-rank approximation & Aproxima la attention global & El error de aproximación y la eficiencia en hardware deben evaluarse en conjunto. \\
\bottomrule
\end{tabularx}
\end{knowledgebox}

\subsection{Del conjunto de datos fijo hacia el meta-learning}

Los experimentos de esta clase corresponden a fixed-dataset supervised learning, pero el input contract se presta de forma natural a una context-based extension: se colocan el support set, el unlabeled set y el query set en una matriz unificada, y la tarea se especifica mediante masks. Si se entrena a través de varios datasets, NPT podría aprender un within-dataset algorithm general; esto lo conecta directamente con in-context learning, few-shot learning, semi-supervised learning y domain adaptation.

\teachervoice{El expositor clasifica explícitamente estas líneas como future work, entre ellas eliminar la mini-batch approximation, la few-shot generalization, domain adaptation, semi-supervised y multi-task learning. Estas notas las presentan solo como líneas de investigación y no describen experimentos aún no realizados como capacidades ya establecidas.}

\subsection{Resumen de la sección}

La capacidad expresiva de la dataset attention y su costo $O(b^2)$ provienen del mismo diseño. GNN, retrieval, sparse attention y meta-learning ofrecen extensiones naturales, pero cada aproximación cambia el context visible y el statistical contract, por lo que requiere validar de nuevo el mecanismo.
